%% notes-tex/031-yeast-chemical-space/sections/2-inventory.tex

\section{What the reconstruction actually names\secstatus{todo}}
\label{sec:inventory}

Figure~\ref{fig:accounting} measures yeast-GEM 9.0.2 as an inventory of chemical species
rather than as a flux model.

%% SOURCE: experiments/031-env-chemgen-inhibitor-tolerance/scripts/yeast9_chemical_accounting.py ->
%% yeast9_accounting.svg, from results/yeast9_accounting_*.csv, converted by `make plots`.
\tcfig{figures/yeast9_accounting.pdf}{\textbf{Yeast9 names 1{,}378
chemical species, and only a fifth of them are pinned to a single molecule.} The reconstruction
read as a chemical inventory rather than as a flux model. \textbf{a}, Distinct species by
structure class, with the share carrying a SMILES. \textbf{b}, Published metabolome counts,
which this script does not measure. \textbf{c}, How precisely a species is named: no structure,
a structure with an unassigned stereocenter that names a set of stereoisomers, or a structure
that names one molecule. \textbf{d}, Share of each reaction's participants carrying a
structure. \textbf{e}, Species reachable from the model's own boundary metabolites when a
reaction fires as soon as any one substrate is present, an upper bound that ignores
stoichiometry and thermodynamics and that saturates in four rounds. \textbf{f}, The species
without a structure, split into those a lookup could recover and those that can never have
one.}{fig:accounting}

Four measurements matter.

\notesec{The inventory is smaller than the model's own headline number}
The release reports 2{,}805 metabolites, but those are compartment-specific entries. Collapsed
to distinct chemical species the model names \textbf{1{,}378}, spread over 14 compartments as
2{,}806 entries in the file we read. Of those species 894 carry a SMILES and 484 do not.

\notesec{Only a fifth of the model is pinned to a single molecule}
Of the 894 structures, \textbf{621 carry at least one unassigned stereocenter}, so the string
names a set of stereoisomers rather than a compound. Only \textbf{273 of 1{,}378 species}, a
fifth of the model, are specified precisely enough for a physical accounting to treat them as
one species. This is the measurement that bears most directly on whether an atom-mapped
treatment is possible, and it is not reported anywhere in the literature we mirror.

\notesec{Half the reactions are structurally complete}
Every participant carries a structure in \textbf{2{,}156 of 4{,}131} reactions, 52 percent.
That is the share a structure-aware or rule-based method could read without first resolving
something, and it is the practical ceiling on applying reaction rules to this model as
released.

\notesec{Connectivity gives no useful bound}
Starting from the model's 267 boundary metabolites and firing a reaction whenever any one of
its substrates is present, \textbf{1{,}374 of 1{,}378 species are reached within four rounds}.
The rule is deliberately generous, and the result is that graph connectivity says only that
the network is dense. It follows that bounding the space of intermediates needs real reaction
rules with atom mapping rather than reachability on a species graph, which is an argument for
the retrobiosynthetic framing rather than against it.

One number to treat carefully. Counting metabolite entries that participate in a single
non-boundary reaction or that can only be produced or only consumed gives \textbf{747 of
2{,}806} on 9.0.2. The published counterpart for Yeast8 is 464 of 2{,}742 \external{Chen, Li and
Nielsen 2022}. The two are not directly comparable because the definitions of a dead end may
differ; ours is stated in the script and is the more inclusive reading.
